\section{CM arithmetic and the classification}\label{sec:enumeration}

The CM classification rests on two structures.  A modular correspondence
bounds the orbit of the elliptic invariant; the complementary CM
eigenline determines the coefficient slope.  The first structure makes
the parameter set finite, and the second produces an identity at every
admissible CM parameter.

\subsection{The modular correspondence}

Write $j_s(x)=j(E_{s,x})$ for the elliptic families of
Section~\ref{sec:families}.  Their invariants are
\begin{equation}\label{eq:jmaps}
\begin{array}{c|c}
\ell&j_{s_\ell}(x)\\\hline
1&1728/[4x(1-x)]\\
2&64(1+3x)^3/[x(1-x)^2]\\
3&27(1+8x)^3/[x(1-x)^3]\\
4&256(1-x+x^2)^3/[x^2(1-x)^2].
\end{array}
\end{equation}
The involution $\iota:x\mapsto1-x$ fixes
$C=R_\ell/[4x(1-x)]$.  At levels one and four it fixes $j_s$ as
well.  At levels two and three it exchanges the invariants of a
cyclically isogenous pair.  Indeed, with
\[
 t_n=\left(\frac{\eta(\tau)}{\eta(n\tau)}\right)^{24/(n-1)},
 \qquad (n,k_n)=(2,64),(3,27),
\]
the classical $X_0(n)$ coordinates give
\[
 x=\frac{k_n}{t_n+k_n},\qquad
 C=\frac{(t_n+k_n)^2}{t_n},\qquad
 w_n(t_n)=\frac{k_n^2}{t_n}.
\]
Here $w_n$ is the Fricke involution, and the two maps to the
$j$-line are
\[
\begin{array}{c|cc}
 n&j(\tau)&j(n\tau)\\\hline
 2&(t_2+256)^3/t_2^2&(t_2+16)^3/t_2\\
 3&(t_3+27)(t_3+243)^3/t_3^3&(t_3+27)(t_3+3)^3/t_3.
\end{array}
\]
Substitution identifies them with $j_{s_n}(x)$ and
$j_{s_n}(1-x)$, respectively.  Thus the
arithmetic information belongs to the correspondence
\[
\begin{tikzcd}[column sep=large]
 X_0(n) \arrow[r,"j"] \arrow[d,"/\langle w_n\rangle"']
   & X(1)\\
 X_0(n)^+ \arrow[r,"C"'] & \mathbf P^1_C .
\end{tikzcd}
\]
There is no asserted map from the lower row to $X(1)$: the missing
descent is precisely why a quadratic equation in $j$ occurs.

The equations of the four correspondences may be written compactly as
\begin{align}
 \Phi_1(J,C)&=C-J,\notag\\
 \Phi_2(J,C)&=J^2-(C^2-207C+3456)J+(C+144)^3,
       \label{eq:P2}\\
 \Phi_3(J,C)&=J^2-C(C^2-126C+2944)J+C(C+192)^3,
       \label{eq:P3}\\
 \Phi_4(J,C)&=(C-16)^3-CJ.\notag
\end{align}
Each is monic as a polynomial in $C$.  For $C\ne0$, its degree in
$J$ is
\[
                         m_\ell=(1,2,2,1)_\ell.
\]
These two elementary features encode, respectively, integrality and
the CM degree bound.

\begin{proposition}[Arithmetic of CM parameters]
\label{prop:rational-cm-integrality}
Let $C\in\Qbar^\times$ have degree at most $d$, and suppose that
the indicated nonsingular elliptic fiber has CM by the order of
discriminant $D$.  Then $C$ is an algebraic integer and
\[
                             h(D)\le m_\ell d.
\]
In particular, a rational CM parameter is an integer; it has
$h(D)\le1$ at levels one and four and $h(D)\le2$ at levels two
and three.  At level four a rational CM parameter also satisfies
$C\mid4096$.
\end{proposition}
\begin{proof}
We use the following precise consequence of the main theorem of
complex multiplication, including nonmaximal orders
\cite[Chapters~7 and~11]{CoxCM}.  If
$\operatorname{End}(E)=\mathcal O_D$, then $j(E)$ is integral and
\begin{equation}\label{eq:cm-main-packet}
 H_D(J)=\prod_{[\mathfrak a]\in\operatorname{Pic}(\mathcal O_D)}
       \bigl(J-j(\mathbf C/\mathfrak a)\bigr)\in\Z[J]
 \quad\text{is irreducible},\qquad
 [\Q(j(E)):\Q]=h(D).
\end{equation}
Artin reciprocity acts simply transitively on these $h(D)$
invariants over the quadratic field; complex conjugation permutes
the same orbit.  The degree of a single invariant over $\Q$ is
therefore $h(D)$, whereas the ring class field is obtained after
adjoining the quadratic field.

Now $\Phi_\ell(j,C)=0$ gives
\[
 h(D)=[\Q(j):\Q]
       \le[\Q(C,j):\Q]\le m_\ell[\Q(C):\Q].
\]
Viewed in the other variable, $\Phi_\ell(j,C)$ is monic over the
ring of algebraic integers.  Transitivity of integrality makes
$C$ integral, and a rational algebraic integer belongs to $\Z$.
For level four the equation also gives
\[
              4096=C\bigl(C^2-48C+768-j\bigr),
\]
whose second factor is an integer when $C$ is rational.
\end{proof}

\subsection{The canonical CM coefficient line}
\label{sec:cm-existence-certificates}

The existence of the coefficient slope is most naturally expressed
in de Rham cohomology.  At a smooth CM fiber over a number field
$K$, put
\[
 H=H^1_{\mathrm{dR}}(E/K),\qquad L=F^1H.
\]
The CM algebra acts on $H\otimes_K\Qbar$ with two eigenlines.
One is $L\otimes_K\Qbar$.  Galois preserves their unordered pair
and fixes $L$, so its companion is also defined over $K$:
\begin{equation}\label{eq:cm-canonical-splitting}
                             H=L\oplus L^{\mathrm{CM}}.
\end{equation}
This splitting is intrinsic and isogeny invariant.  It does not
require all geometric endomorphisms to be defined over $K$.

In the standard models write $\Omega=dX/y$ and $\Xi=X\,dX/y$,
and put $H_x=H^1_{\mathrm{dR}}(E_{s,x}/K)$.  For
$x\ne0,1,1/2$, define
\[
 Q(x)=\frac{2x(1-x)}{1-2x},\qquad
 T_x:K^2\longrightarrow H_x,\quad
 (A,B)\longmapsto A\Omega_x+BQ(x)\nabla_x\Omega_x.
\]
The Gauss--Manin formula \eqref{eq:logderivative} shows that
$T_x$ is an isomorphism: in the basis $(\Omega,\Xi)$ its determinant
is $v_s(x)\ne0$.  The CM coefficient line is consequently
\begin{equation}\label{eq:cm-coefficient-line}
                       \mathcal R_x=T_x^{-1}(L_x^{\mathrm{CM}}).
\end{equation}
It is a $K$-line, and it is different from $B=0$.

\begin{proposition}[The endomorphism determinant]
\label{prop:cm-matrix-certificate}
Let $z\in[-1,1)$ be nonzero and algebraic, and suppose that its indicated
elliptic fiber has CM.  A nonzero coefficient vector $(A,B)$ on
$\mathcal R_x$ satisfies
\[
                 (A+B\thetaop)F_s(z)=\frac\alpha\pi,
                 \qquad \alpha\in\Qbar^\times.
\]
More explicitly, choose a nonrational endomorphism $\varphi$ and
write
\[
 \varphi^*\Omega=\lambda\Omega,\qquad
 \varphi^*\Xi=r\Omega+\bar\lambda\Xi,
 \qquad \kappa=\frac{r}{\bar\lambda-\lambda}.
\]
If $d$ is the signed index of
$(\gamma,\varphi_*\gamma)$ in an oriented integral symplectic
homology basis, then
\begin{align}
 \frac AB&=v\kappa-u,\label{eq:cm-matrix-slope}\\
 \alpha&=\frac{3iBvd}{\bar\lambda-\lambda}.
 \label{eq:cm-matrix-alpha}
\end{align}
Here $u,v$ are those of \eqref{eq:logderivative}, and $d\ne0$.
\end{proposition}
\begin{proof}
The vector $\nu=\kappa\Omega+\Xi$ spans the complementary
eigenline.  Write $\omega=\int_\gamma\Omega$ and
$\rho=\int_\gamma\nu$.  The intrinsic determinant identity
\eqref{eq:general-cm-determinant} gives $d\ne0$ and
\[
                  (\bar\lambda-\lambda)\omega\rho=2\pi i d.
\]
The condition $T_x(A,B)\in\Qbar\nu$ says
$A+Bu=Bv\kappa$.  The normalization of the standard period now
gives
\[
 (A+B\thetaop)F_s
   =\frac{3Bv}{2\pi^2}\omega\rho
   =\frac1\pi\frac{3iBvd}{\bar\lambda-\lambda}.
\]
The multiplier is algebraic and nonzero.  The identity-line theorem
identifies this line with the unique possible algebraic coefficient
slope in the ordinary convergence interval.
\end{proof}

\begin{lemma}[Equivariance of the coefficient line]
\label{lem:cm-involution-covariance}
The parameter involution $x\mapsto1-x$ identifies projective
coefficient vectors under an isogeny of the corresponding smooth
elliptic fibers.  In particular, it preserves the line
$\mathcal R_x$ whenever the fiber has CM.
\end{lemma}
\begin{proof}
The modular correspondence supplies an isogeny
$\psi:E_{s,x}\to E_{s,1-x}$ of degree $1,2,3,1$, respectively,
over an algebraic extension of the function field.  For levels one
and four the equality of the $j$-invariants supplies an isomorphism;
for the middle levels this is the cyclic isogeny on $X_0(2)$ or
$X_0(3)$ just described.  Write
$\psi^*\Omega_{1-x}=c(x)\Omega_x$.

The multiplier is constant.  Indeed, the period equations for the
two chosen differentials are the same Gauss equation
\[
 y''+\frac{1-2x}{x(1-x)}y'
                 -\frac{s(1-s)}{x(1-x)}y=0,
\]
which is invariant under $x\mapsto1-x$.  On a simply connected
open set, $\psi$ is a constant map of Betti local systems.
Multiplication of a solution basis by $c(x)$ multiplies its
Wronskian by $c(x)^2$; both Wronskians are constant multiples of
$[x(1-x)]^{-1}$.  Thus $c(x)=c\in\Qbar^\times$.

Functoriality of Gauss--Manin and $Q(1-x)=-Q(x)$ give
$\psi^*T_{1-x}=cT_x$.  After projectivizing, the square
\begin{equation}\label{eq:coefficient-covariance}
\begin{tikzcd}[column sep=large]
 \mathbf P^1(\Qbar) \arrow[r,"\mathbf P(T_{1-x})"]
                        \arrow[d,equal]
   & \mathbf P(H_{1-x}\otimes\Qbar)
                        \arrow[d,"\mathbf P(\psi^*)"]\\
 \mathbf P^1(\Qbar) \arrow[r,"\mathbf P(T_x)"']
   & \mathbf P(H_x\otimes\Qbar)
\end{tikzcd}
\end{equation}
commutes.  An isogeny intertwines the CM endomorphism algebras and
preserves the Hodge lines, hence also their companions.  This proves
the assertion.  The identity extends to every smooth fiber with
$x\ne1/2$ by extension of homomorphisms of abelian schemes; the
constant nonzero differential multiplier remains nonzero.
\end{proof}

\begin{theorem}[Rational descent of CM identities]
\label{thm:rational-cm-existence}
Let $z\in\Q\setminus\{0\}$ with $-1\le z<1$.  If its indicated
elliptic fiber has CM, there is a primitive integer pair $(A,B)$,
unique up to sign, for which the ordinarily convergent standard
series is a nonzero algebraic multiple of $1/\pi$.
\end{theorem}
\begin{proof}
Set $K=\Q(x)$, where $x=(1-\sqrt{1-z})/2$.  The line
$\mathcal R_x$ is defined over $K$ by
\eqref{eq:cm-canonical-splitting}.  If $K=\Q$ there is nothing
to descend.  Otherwise its nontrivial automorphism $\sigma$ sends
$x$ to $1-x$.  Galois functoriality gives
$\mathcal R_x^\sigma=\mathcal R_{1-x}$, while the commutative
square \eqref{eq:coefficient-covariance} gives
$\mathcal R_{1-x}=\mathcal R_x$.  Therefore
$\mathcal R_x\in\mathbf P^1(\Q)$.

Choose its relatively prime integer representative.  The
endomorphism determinant gives the nonzero algebraic multiplier,
and Theorem~\ref{thm:unconditional} gives uniqueness.  The
convergence statement follows from the assumed interval, including
ordinary alternating convergence at $z=-1$.
\end{proof}

\subsection{Finite exhaustion}

Only orders of class number at most two remain.  To make this
reduction effective, we use the proved classification of quadratic
fields in \cite{Watkins}: a negative fundamental discriminant with
class number at most two satisfies $|d|\le427$.  For an arbitrary
order of discriminant $D=f^2d$, the conductor formula is
\begin{equation}\label{eq:conductor-packet}
 h(f^2d)=\frac{h(d)f}{[\mathcal O_d^\times:
                              \mathcal O_{f^2d}^\times]}
            \prod_{p\mid f}\left(1-\frac{(d/p)}p\right).
\end{equation}
The order class group surjects onto the maximal-order class group,
so $h(d)\le h(f^2d)$.  The unit index is at most three and
\begin{equation}\label{eq:conductor-bound-packet}
 h(f^2d)\ge\frac{h(d)}3\varphi(f)
                  \ge\frac{h(d)}3\sqrt{f/2}.
\end{equation}
The last inequality follows by multiplying
$p^{a-2}(p-1)^2\ge1$ over odd prime powers $p^a\Vert f$;
the factor at two is at least $1/2$.  Thus $h(D)\le2$ implies
$f\le72$.  The finite range is a theorem, not a numerical cutoff.

\begin{lemma}[Exact arithmetic enumeration]
\label{lem:cm-finite-enumeration}
The rational CM parameters in the closed convergence domain
$|C|\ge R_\ell$ are the $C$-columns of
Tables~\ref{tab:standard1}--\ref{tab:standard4}, together with
$C=R_\ell$ at each level.
\end{lemma}
\begin{proof}
Primitive reduced positive definite binary quadratic forms and
\eqref{eq:conductor-packet}, in the proved range
$-427\le d\le-3$, $1\le f\le72$, give thirteen orders of class
number one and twenty-nine of class number two.  For each order
compute its integral class polynomial $H_D$.

There is a short elimination rule.  If $h(D)=1$, substitute its
single integral root in $\Phi_\ell(J,C)$.  If $h(D)=2$, only the
middle levels are possible; a shared root forces the monic
quadratics $H_D(J)$ and $\Phi_\ell(J,C)$ to be equal.  Thus, writing
them as $J^2-T_DJ+N_D$ and $J^2-T_\ell(C)J+N_\ell(C)$, the
possible $C$ are exactly the integer roots of
\[
               \gcd\bigl(T_\ell(C)-T_D,
                          N_\ell(C)-N_D\bigr).
\]
Exact factorization therefore lists every possible parameter
without a bound on its magnitude.  It gives
\[
\begin{array}{c|rrrr}
\ell&1&2&3&4\\\hline
|C|\ge R_\ell&12&13&10&5\\
C\ne R_\ell&11&12&9&4.
\end{array}
\]
The complete discriminants, class polynomials, and factorizations
are retained in the finite arithmetic certificate described in
Appendix~\ref{app:cm-certificate}.  The class
polynomial coefficients were certified by complex interval
arithmetic and their integrality; independent PARI calculations
gave the same polynomials.  Integer roots were recovered both by
factorization and, for the middle levels, by elimination resultants.

The correspondence may initially retain either elliptic invariant
above $C$.  They are isogenous, so CM of one is equivalent to CM of
the other.  Every retained nonsingular fiber is consequently CM on
the specified branch.  This proves both directions of the
enumeration.
\end{proof}

\begin{theorem}[The standard CM classification]\label{thm:conditional}
The primitive ordinarily convergent standard identities whose
indicated elliptic fibers have CM are exactly the 36 rows of
Tables~\ref{tab:standard1}--\ref{tab:standard4}, up to simultaneous
sign.  Under Hypothesis~\ref{hyp:QPI}, they exhaust all primitive
standard identities.  The CM assertion remains unchanged if the
nonzero parameter $C$ is allowed to be rational.
\end{theorem}
\begin{proof}
Integrality reduces rational CM parameters to integer ones.
Lemma~\ref{lem:cm-finite-enumeration} exhausts their possible values
in the convergence domain.  Theorem~\ref{thm:unconditional}
excludes the four positive endpoints, and rational descent gives
an identity at every remaining parameter.  It also permits only
one primitive pair up to sign.  The evaluations below identify
those pairs and their multipliers with the tabulated rows.
Finally, conditional CM forcing supplies the assertion without
the CM premise.  The conditional assertion uses the parameter domain of
Hypothesis~\ref{hyp:QPI}; enlarging that domain to arbitrary
rational $C$ requires enlarging the hypothesis as well.
\end{proof}

\subsection{Exact evaluation of the classified lines}
\label{sec:standard}

The geometric proof establishes existence and rationality before
computing a single singular value.  For the displayed coefficients
we use the exact modular evaluations in
\cite[Tables~5--11]{CCY}.  Their role is specific: with the
normalizations
\[
 X_\ell(q)=q+O(q^2),\qquad
 Z_\ell(q)=\sum_{n\ge0}u_\ell(n)X_\ell(q)^n,\qquad
 G_\ell(X)=1-R_\ell X,
\]
these tables give the singular value $X_\ell=1/C$ and the
coefficient slope $\lambda=A/B$ obtained by differentiating the
corresponding modular equation.  The modular summation theorem of
Chan--Chan--Liu \cite[Theorem~2.1]{CCL}, in this normalization
\cite[Theorem~2.4]{CCY}, applies to the specified weight-two form
on the indicated normal extension of $\Gamma_0(\ell)$.  At the
tabulated quadratic nome, inside the open convergence disk, its
conclusion is
\begin{equation}\label{eq:alpha}
 \alpha=\frac{B}{2\sqrt{1-R_\ell/C}}
 \begin{cases}
   \sqrt{\ell/N},&q=e^{-2\pi\sqrt{N/\ell}},\\
   \sqrt{L/N},&q=-e^{-2\pi\sqrt{N/L}},
 \end{cases}
 \qquad L=\frac{4\ell}{\gcd(4,\ell)}.
\end{equation}
Thus each interior row follows from two exact modular values and
this summation theorem; its displayed radical is then fixed by
\eqref{eq:alpha}.  These evaluations furnish identities on the
CM locus, not a converse for arbitrary algebraic-value identities.

The remaining row is the boundary identity of Bauer,
\[
 \sum_{n\ge0}(4n+1)\frac{(1/2)_n^3}{(n!)^3}(-1)^n
       =\frac2\pi,
\]
as recorded with its original source in
\cite[Section~4.1 and Table~11]{CCY}.  Its ordinary convergence
was proved in the analytic part.  It is the sole conditionally
convergent row; all other rows are absolutely convergent.

\subsection{Linear Euler companions}\label{sec:euler}

Put
\[
 G_s(z)={}_2F_1(1-s/2,(1+s)/2;1;z)^2.
\]
Euler's transformation identifies the two branches and their
coefficient lines:
\begin{equation}\label{eq:euler}
 F_s=(1-z)G_s,\qquad
 (A+B\thetaop)F_s
   =\bigl(A(1-z)-Bz+B(1-z)\thetaop\bigr)G_s.
\end{equation}
This is a change of differential basis, whose determinant is
$(1-z)^2$.  Away from $z=1$ it induces an isomorphism of rational
projective coefficient lines.

\begin{theorem}[Euler transport]\label{thm:euler}
For $|C|>R=R_\ell$, define
\[
 (a_0,b_0)=\bigl((C-R)A-RB,(C-R)B\bigr),\qquad
 d=\gcd(|a_0|,|b_0|).
\]
The primitive Euler companion of a primitive standard identity is
\[
                  (a,b,\beta)=\varepsilon
                  (a_0/d,b_0/d,C\alpha/d),
\]
where $\varepsilon=\pm1$ chooses the representative.  This
transport is bijective on rational projective slopes.  The
convergent CM companions are exactly the 35 rows of
Tables~\ref{tab:euler1}--\ref{tab:euler4}; under
Hypothesis~\ref{hyp:QPI}, they are all the primitive companions.
No nonzero companion pair converges at either endpoint.
\end{theorem}
\begin{proof}
Multiply \eqref{eq:euler} by $C$.  The coefficient matrix is
\[
 M_C=\begin{pmatrix}C-R&-R\\0&C-R\end{pmatrix},
 \qquad \det M_C=(C-R)^2.
\]
Its inverse on projective lines is represented by
$(a,b)\mapsto((C-R)a+Rb,(C-R)b)$.  Both directions preserve
algebraicity of the multiplier, and division by the gcd gives
exactly the primitive representative.

For convergence, the coefficients of $G_s=F_s/(1-z)$ are
$d_s(n)=\sum_{k=0}^n c_s(k)$.  Positivity and convergence of
$F_s(1)$ give $d_s(n)\to F_s(1)>0$.  Hence
$(a+bn)d_s(n)z^n$ fails to tend to zero for every nonzero pair
when $|z|\ge1$.  In $|z|<1$ the series converges absolutely.
Thus transport preserves all 35 interior standard rows and loses
the single Bauer boundary row.  The standard classification and
the inverse transport prove exhaustion, with exactly the same
conditional qualification.
\end{proof}
